\scanpage{141}
\mspage{70}

L'image de cette application est donc formée des \inserted{classes de $\pi$-conj. des} sous-groupes sections $\Gamma'\subset\mathcal E$ tels que $\widetilde U^{\Gamma'}\ne\varnothing$, et pour un tel $\Gamma'$, l'ens. des \unsure{$x\in U^\Gamma$} qui \deleted{\unsure{lui}} donnent comme image cette classe de conj. est l'image de $U^{\Gamma'}$ dans $U$. Enfin, si $x$ est dans cette image, l'ens. \inserted{$\widetilde U^{\Gamma'}_x$} des $\widetilde x\in\widetilde U^{\Gamma'}$ au dessus de $x$ (\deleted{\unsure{$\ne\varnothing$}} hyp. sur \unsure{a|x}) est un torseur sous $\pi^{\Gamma'}$ i.e. si $\widetilde x\in(\widetilde U^{\Gamma'})_x$, \deleted{\unsure{de}} et si $\alpha\in\pi$, alors
\[
\alpha.\widetilde x\in\widetilde U^{\Gamma'}\quad\Longleftrightarrow\quad\alpha\in\pi^{\Gamma'}
\]
(Vérification triviale, comme du reste toutes les assertions précédentes). Ainsi, les conditions suivantes sont équivalentes :

\margin{NB \unsure{$\widetilde U_x$} étant identifié à \unsure{$\operatorname{Isom}_U(R_x,\widetilde U)$}, $\widetilde U_x^{\Gamma'}$ s'identifie à \unsure{$\operatorname{Isom}_{U,\Gamma}(R_x,\widetilde U_{\Gamma'})$} \unsure{des deux isom. qui commutent|des isom. de rev. commutant} à l'action de $\Gamma$ et $\pi^\Gamma=(1)$ signifie que cet isom. est unique \unsure{i.e.} \deleted{\illegible{3}} \unsure{la classe de conj. des sections} \illegible{2}.}

 a) $(*)$ est une bijection, et $\forall$ ss-gpe $\Gamma'$ de $\mathcal E$ section de l'extension, \deleted{\illegible{1}} on a $\pi^{\Gamma'}=\{1\}$.

 b) $\forall\,\Gamma'$ comme dans a), $\Gamma'$ opérant sur $\widetilde U$ a un p\textsuperscript{t} fixe et un seul.

Ceci posé, prouvons le

\uline{Lemme fondamental} Soit $\Gamma$ un groupe \inserted{fini} fini, opérant fidèlement sur \inserted{un espace} \unsure{$D$}\,$\simeq\mathbb R^2$, soit \deleted{On suppose $\Gamma=$} $\Gamma^\circ$ le ss-gpe de $\Gamma$

\mspage{140}

\scanpage{142}

(d'indice $0$ ou $1$) formé des $g\in\Gamma$ tels que $g_D$ conserve l'orientation, et supposons $\Gamma^\circ\ne\{1\}$ (i.e. $\Gamma$ n'est réduit ni au groupe unité, ni au groupe \deleted{\illegible{1}} $\{1,\sigma\}$, où $\sigma$ est une anti-involution de $D$. Alors

 a) $\Gamma$ admet un p\textsuperscript{t} fixe et un seul dans $D$ i.e. $\operatorname{card}D^\Gamma=1$

 b) \deleted{\unsure{Et}} $\Gamma^\circ$ est cyclique, et si $\Gamma\ne\Gamma^\circ$, $\Gamma$ est un groupe \inserted{diédral}.

Plus précisément, il existe un homéomorphisme
\[
D\simeq\mathbb C
\]
tel que \deleted{$\Gamma$ opère par automorphismes} \inserted{\unsure{sur $\mathbb C$ par le groupe}} \deleted{\illegible{1}} d'homéomorphismes de $\mathbb C$ transformé de \deleted{complexes et anti-complexes de $\mathbb C$} \deleted{\unsure{fixant}} $\Gamma$ soit le groupe des homothéties par $\mu_n(\mathbb C)$ (si $\Gamma=\Gamma^\circ$ d'ordre $n$), soit le groupe diédral associé \unsure{$z\longrightarrow\zeta\tau^\varepsilon(z)$} ($\zeta\in\mu_n(\mathbb C)$, $\varepsilon=\pm1$) où \unsure{$\sigma$} \deleted{\illegible{1}} est la conjugaison complexe.

\uline{Démonstration du lemme fondamental.}

\deleted{Supposons d'abord $\Gamma=\Gamma^\circ$, i.e. $\Gamma$ conserve l'orientation. \deleted{Si $\Gamma$ opérait \unsure{sans p\textsuperscript{t}}} On procède}

\mspage{141}

\scanpage{143}
\mspage{71}

\deleted{Par induction sur l'ordre de $\Gamma$.}

\deleted{\unsure{a)} \unsure{$D^\Gamma\ne\varnothing$}}

\deleted{a) Cas $\Gamma$ d'ordre premier}

 a) Supposons d'abord qu'on puisse trouver une structure $C^2$ sur $D$ invariante par $\Gamma$ --- alors un argument standard montre qu'il existe une structure conforme invariante par $\Gamma$, or le th. fondamental de la représentation conforme montre qu'alors
\[
\begin{aligned}
&\text{ou bien }D\simeq\text{intérieur}\ \inserted{\Delta}\ \text{du disque unité }\Delta\ \inserted{\text{ou du demi-pl. de Poincaré}}\\
&\text{ou bien }D\simeq\mathbb C\qquad\text{(isom. conformes).}
\end{aligned}
\]
Dans le premier cas, le groupe des automorphismes \deleted{applications} conformes est
\[
\simeq\unsure{\operatorname{Sl}(2,\mathbb R)^\sim\,|\,\operatorname{Gl}(2,\mathbb R)^\sim}
=\{u\in\operatorname{Gl}(2,\mathbb R)\mid\det u=\pm1\}
\]
(\deleted{opérant par} $u=\left(\begin{smallmatrix}a&b\\c&d\end{smallmatrix}\right)$ opérant sur \deleted{\illegible{1} Poincaré} \inserted{$\widehat{\mathbb C}$} par \deleted{\unsure{$z\longmapsto\dfrac{az+b}{cz+d}$}} $\Theta_u\circ\tau^{\det u}$, où $\tau$ est la conjugaison complexe et $\Theta_u(z)=\dfrac{az+b}{cz+d}$, et laissant stable le demi-pl. de Poincaré), et tt sous-groupe compact est \inserted{contenu dans un} conjugué du \uline{ss-groupe} compact maximal, qui (en repassant au disque unité \inserted{$\Delta$}) s'identifie au groupe $O(2,\mathbb R)$ des

\mspage{142}

\scanpage{144}

transformations du disque unité de la forme
\[
z\longrightarrow\zeta\tau^\varepsilon(z)\qquad
\begin{gathered}
\zeta\in\mathbb U=\{\zeta\in\mathbb C\mid |\zeta|=1\}\\
\varepsilon\in\{\pm1\}\\
\tau\ \text{la conjugaison complexe}
\end{gathered}
\]
Tout ss-groupe fini de ce groupe est de l'un des types explicités plus haut --- on gagne, en utilisant un \unsure{homéom.} $[0,1[\longrightarrow[0,+\infty[$ pour définir un homéom \deleted{\illegible{1}} $\Delta\simeq\mathbb C$ commutant à l'action de $G=\deleted{\unsure{S}}O(2,\mathbb R)$.

Dans le cas $D\simeq\mathbb C$, on voit que \deleted{\unsure{est}} le groupe des autom. conformes de $\mathbb C$ est le groupe des transformations $az+b$ ou $a\overline z+b$, dans lequel un ss-groupe compact maximal est le $\overline{\mathrm m}$ groupe $\deleted{\unsure{S}}O(2,\mathbb R)$ que tantôt --- et tt ss-groupe compact (a fortiori, tt ss-groupe fini) est contenu dans un conjugué de celui-ci. On gagne encore.

Le reste du travail consiste \inserted{ess\textsuperscript{t}} à \deleted{éliminer} montrer que l'hyp. de non sauvagerie

\mspage{143}

\scanpage{145}
\mspage{72}

est toujours satisfaite, \inserted{du moins pour $\Gamma^\circ$. Supposons d'abord $\Gamma=\Gamma^\circ$.} \deleted{On suppose} \inserted{On suppose} que tt est prouvé pour \deleted{les groupes} \inserted{les ordres} $<\operatorname{card}(\Gamma)$. \deleted{Supposons}

 b) $\Gamma$ admet un p\textsuperscript{t} fixe \inserted{\unsure{$o$|$0$}}, i.e. $D^\Gamma\ne\varnothing$.

Sinon, les ss-groupes de stabilité \inserted{$\Gamma_x$} des $x\in D$ étant d'ordre $<$ \deleted{\illegible{1}} $\operatorname{card}\Gamma$, par hyp. de récurrence, les $\Gamma_x$ ont la structure dite dans le th., donc \unsure{$D^\Gamma=U$|$D/\Gamma=U$} est une surface topologique et \unsure{$D\longrightarrow U$|$D\longrightarrow D^\Gamma$} est un rev ramifié, choisissons une structure \deleted{conforme complexe} \inserted{conforme} sur $U$ \deleted{NB $U$ est orienté}; il y a une unique structure conforme sur $D$ telle que $D\longrightarrow U$ soit \deleted{\unsure{conforme|holomorphe}} \inserted{``conforme''} (hol. ou anti-holomorphe), celle-ci est invariante par $\Gamma$, et d'après a) $\Gamma$ admet un p\textsuperscript{t} fixe, contradiction.

 c) $\Gamma$ n'admet pas d'autre p\textsuperscript{t} fixe que \unsure{$o$|$0$}. On fait opérer $\Gamma$ \inserted{\unsure{fidèl\textsuperscript{t}}} sur \unsure{$D^*=D\setminus\{o\}$|$D^*=D\setminus\{0\}$}, et il faut prouver que $D^{*\Gamma}=\varnothing$.

\mspage{144}

\scanpage{146}

Soit donc $x\in D^{*\Gamma}$ --- on va aboutir à une contradiction. Considérons le rev. universel $\widetilde D^*$ de $D^*$ ponctué en $x$, \deleted{OPS} donc $\Gamma$ opère sur $\widetilde D^*$ \unsure{avec} p\textsuperscript{t} fixe $\widetilde x$ au dessus de $x$. Ici $\pi=\pi_1(D^*)\simeq\mathbb Z$, et $\Gamma$ y opère \unsure{trivial\textsuperscript{t}} (car $\Gamma=\Gamma^\circ$) donc $\Gamma\times\mathbb Z$ opère sur $\widetilde D^*$. OPS $D=\mathbb C$, \unsure{$o=0$|$0=0$}, $x=1$, $D^*=\mathbb C^*$, $\widetilde D^*=\mathbb C$, $\widetilde x=0$, $\widetilde D^*\longrightarrow\widetilde D$ donné par $\exp 2i\pi z$, \deleted{\unsure{donc}} et $\mathbb Z$ opérant sur $\mathbb C$ par $\Theta_nz=z+n$ ($n\in\mathbb Z$).

\deleted{Lemme} Il reste à prouver que si un \inserted{groupe fini} $\Gamma$ opère sur $\mathbb C$ en commutant \deleted{\unsure{à}} à l'action de $\mathbb Z$ sur $\mathbb C$, et en laissant \deleted{un p\textsuperscript{t}} fixe le p\textsuperscript{t} \unsure{$o$|$0$}, alors $\Gamma$ opère trivialement (ce qui contredit l'hyp. de fidélité de l'opération).

\deleted{Donc} On est ramené au.

\uline{Lemme} Soit $u$ un homéomorphisme \inserted{d'ordre fini} de $\mathbb C$,

\mspage{145}

\scanpage{147}
\mspage{73}

commutant à $z\longmapsto z+1$, et laissant invariant l'origine. Alors $u=\operatorname{id}$

OPS \deleted{$u$ d'ordre premier} $\exists$ nb premier $p$ tel que $u^p=\operatorname{id}$, i.e. que $u$ correspond à une opération de $\Gamma=\mathbb Z/p\mathbb Z$ sur $\mathbb C$. \inserted{Tous les p\textsuperscript{ts} de $\mathbb Z\subset\mathbb C$ sont fixes par $\Gamma$.} Passant à $\widehat{\mathbb C}\simeq S^2$, on trouve que $\infty$ est un p\textsuperscript{t} d'accumulation des p\textsuperscript{ts} fixes sous $\Gamma$ --- si $\Gamma$ n'opérait pas trivialement, ce serait décidément très sauvage ! On doit pouvoir terminer par la suite spectrale d'Adams... je n'entre pas dans ces dédales...

 d) La partie purement top. étant \unsure{admise|ainsi} \unsure{supposée|supposé} \unsure{prouvée|prouvé}, on en conclut ainsi, si $\Gamma\ne\Gamma^\circ$, $\Gamma^\circ\ne(1)$, \deleted{que l'unique} \inserted{comme} \deleted{p\textsuperscript{t} fixe de} \inserted{$D^{\Gamma^\circ}$} est invariant sous $\Gamma$, \deleted{que} \inserted{comme} $D^{\Gamma^\circ}$ est réduit à un p\textsuperscript{t}, celui-ci est invariant sous $\Gamma$ tt entier, pas seulement $\Gamma^\circ$. D'autre part, On en sait assez maintenant pour savoir que si $\Gamma$ groupe fini opère sur une surface \deleted{séparée} \inserted{paracompacte $U$}, les $\Gamma_x$ ($x\in U$) respectant l'orientation, alors \unsure{$U/\Gamma=V$}

\mspage{146}

\scanpage{148}

est une surface --- \inserted{et $U\longrightarrow V=U/\Gamma$ est un rev. ramifié} et choisissant une structure conforme sur celle-ci, on trouve par image inverse une str. conforme sur $U$ invariante par $\Gamma$, \deleted{et} \inserted{Pour le cas $U=D$} on termine par a), \deleted{mais} pour le complément du lemme fondamental. [Mais pour bien faire, il faudrait prouver qu'il existe tjrs structure conforme invariante, \deleted{\illegible{2}} $\Gamma$ gpe fini opère sur surfaces \deleted{\illegible{1}} \unsure{paracompactes} --- donc $U/\Gamma$ est une surface à bord.... Ici ce qui manque, c'est l'analyse de l'action d'une anti-involution d'une surface au voisinage d'un p\textsuperscript{t} fixe...]

\margin{Il faudrait prouver que si $\tau$ est un \unsure{anti-auto.} involutif de $D$, alors $\exists$ isom $D\simeq\mathbb C$ tel que $\tau$ devienne $z\longmapsto\overline z$ (donc $D^\tau\simeq\mathbb R$ !) ce qui doit permettre de prouver, si $\Gamma\simeq\mathbb Z/2\mathbb Z$ opère par anti-autom. sur $U$ (\unsure{orient.|orientable} $U\ne S^2$) que $\pi_0(U^\Gamma)$ \unsure{$\Longrightarrow$} \unsure{Classes de $\pi$-conj. de sections de $\mathcal E$ sur $\Gamma$} est bij.}

Conséquence du lemme fondamental :

\uline{Théorème.} Soit $U$ une surface topologique paracompacte \inserted{$0$-connexe}, $\Gamma$ un groupe fini opérant \uline{fidèlement} sur $U$, on suppose $\widetilde U$ \deleted{\illegible{1}} non compact (i.e. $U$ non homéomorphe à $S^2$ ni

\mspage{147}

\scanpage{149}
\mspage{74}

au plan projectif réel), on suppose que \deleted{\unsure{$\Gamma\ne(1)$}, \illegible{1}} que \inserted{de plus} si $U$ est \deleted{non} orientable, alors \deleted{\unsure{$\Gamma\ne\{\pm1\}$} i.e. $\operatorname{card}\Gamma\ge3$ \illegible{1}} \deleted{on} suppose que le sous-groupe $\Gamma^\circ$ de $\Gamma$ des $g\in\Gamma$ qui conservent une orientation soit $\ne(1)$ [donc $\Gamma$ \unsure{non réduit} ni à $1$, ni à \unsure{$1$ et|Id} une anti-involution], et si $U$ non orientable, que \deleted{\unsure{card}} $\Gamma$ \unsure{$\not\simeq$|$\ne$} $(1),\{\pm1\}$ i.e. $\operatorname{card}\Gamma\ge3$. Ceci posé, \deleted{\unsure{l'app.}} considérons \deleted{l'application} \inserted{l'extension} $\mathcal E$ de $\Gamma$ par $\pi=\pi_1(U)$, et l'application
\[
U^\Gamma\longrightarrow\begin{gathered}
\text{Classes de }\pi\text{-conjugaison de}\\
\text{scindages de }\mathcal E\longrightarrow\Gamma
\end{gathered}
\]
on a ceci :

 a) Cette application est bijective

 b) Pour tt ss-groupe section $\Gamma'\subset\mathcal E$, on a $\pi^{\Gamma'}=(1)$

 c) Si $U^\Gamma\ne\varnothing$ i.e. S'il existe un scindage, alors $\Gamma$ est cyclique ou diédral (et dans le cas $U$ orienté, $\Gamma^\circ$ est cyclique)

\uline{NB} Les hypothèses sur $U$ assurent que $\widetilde U\simeq D$, et celles sur $\Gamma$ que $\Gamma$ opérant sur $D$ satisfait aux hyp. du lemme fondamental.

\mspage{148}

\scanpage{150}

\deleted{\uline{Corollaire-Scholie}, Pour un $x\in U^\Gamma$ correspondant à une classe de scindage de l'extension $\mathcal E$, on trouve que les scindages de cette classe, \deleted{sur lesquels $\mathcal E$ opère par \unsure{autom. int.|fonctorialité}}, forment un \deleted{\unsure{$\mathcal E$-ensemble qui est un}} $\pi$-torseur, \deleted{et les automorphismes de}}

\uline{Corollaire} Supposons \deleted{que} $U$ \unsure{orientable}, \deleted{\illegible{1}} $\Gamma$ conservant l'orientation. Soit $U^!$ l'ens des $x\in U$ tels que $\Gamma_x\ne(1)$. (qui est donc une partie discrète de $U$).

À tt $x\in U^!$, associons la classe de conjugaison des ss-groupes de $\mathcal E$ \inserted{sections \inserted{\unsure{$\pi$}} partielles de $\mathcal E$ sur $\Gamma_x$} qui \unsure{correspondent|correspondant} à \unsure{cet|ce} $x\in U^{\Gamma_x}$. Alors

 a) \deleted{l'application} les sections partielles ainsi obtenues sont \uline{maximales} parmi celles qui sont $\ne(1)$

 b) L'application de $U^!$ vers \unsure{l'ens|l'un} des classes de $\pi$-conjugaison des sections partielles \deleted{non triviales} \inserted{$\ne(1)$} maximales, est bijective

\margin{\uline{NB} Cette application commute aux actions naturelles de $\Gamma$ !}

 c) Pour tt telle section partielle, on a $\pi^{\Gamma'}=(1)$, i.e. les \deleted{\unsure{$\Gamma_x$-}} \unsure{automorphismes} du rev. \inserted{universel} \unsure{$R_x$} de $U$ qui commutent à l'action de $\Gamma_x$ sont triviaux

--- il y a un isom. unique \inserted{commutant à l'action de $\Gamma_x$} entre ce torseur, et le torseur déduit de $\widetilde U$ en tordant par le $\pi$-torseur \deleted{\unsure{des}} $P_x$ des \unsure{$\Gamma''$} dans la classe \deleted{\illegible{1}} \unsure{$[\Gamma']$}...

\mspage{149}
